%% parts/p06/large-cardinals-absoluteness.tex --- Phase-4 P6
%% hierarchy/reflection/large-cardinal depth slice.
%%
%% Register ownership: D-P7-006/007/008, N-a-044/045/067, A-065.
%% A-130 is already realized by fig:p06-cumulative-hierarchy; the matching
%% master consistency table/figure follow this chapter in their own live file.
%%
%% Source accounting: v1.1 Part VII contains the cumulative-hierarchy and L
%% developments but no dedicated Levy-hierarchy, Shoenfield-absoluteness, or
%% large-cardinal entry.  Its hierarchy sketch remains
%% thm:p06-hierarchy -> annex:p6.hierarchy.1; its L sketch remains
%% thm:p06-constructible-class -> annex:p6.force.1.  This file adds cited
%% register-driven material and does not promote either source status.

\chapter{Reflection, Absoluteness, and Large Cardinals}
\label{sec:p06-large-cardinals-absoluteness}
\chaptermark{Absoluteness and Large Cardinals}

This chapter supplies the syntactic complexity scale used by absoluteness
arguments and the conditional large-cardinal ladder used to calibrate the
countable-transitive-model convention.  It assumes the set-theoretic language
and the ambient \(\mathsf{ZFC}\) conventions fixed earlier in Part~6.  The
standard references are \cite{Jech2003,Kanamori2003,Moschovakis2009}; Scott's
constructibility obstruction is cited to its primary source
\cite{Scott1961}.

\section{The L\'evy hierarchy and elementary absoluteness}
\label{sec:p06-levy-hierarchy}

\begin{hbdefinition}{def:p06-levy-hierarchy}{L\'evy hierarchy}
A membership-language quantifier is \emph{bounded} when it has one of the
forms
\[
  \exists x\in y\,\varphi
  \quad\text{or}\quad
  \forall x\in y\,\varphi,
\]
which abbreviate quantification restricted by \(x\in y\).  The classes
\(\Delta_0=\Sigma_0=\Pi_0\) consist of formulas all of whose quantifiers are
bounded.  Inductively, the normal-form representatives of
\(\Sigma_{n+1}\) begin with a finite block of unbounded existential
quantifiers followed by a \(\Pi_n\) representative; \(\Pi_{n+1}\) is the
dual class beginning with universal quantifiers.  The representative classes
are closed under bounded quantification and finite conjunction and
disjunction, while negation interchanges \(\Sigma_n\) and \(\Pi_n\).  For
\(n\geq1\), an arbitrary formula is classified as \(\Sigma_n\) or \(\Pi_n\)
relative to the declared set-theoretic base when that base proves it
equivalent to a representative of the indicated class.  For \(n\geq1\), a
formula is \(\Delta_n\) when the base proves it equivalent both to a \(\Sigma_n\)
representative and to a \(\Pi_n\) representative.  The symbol \(\Delta_0\)
continues to mean the literal bounded-quantifier class fixed above.

The hierarchy is metalinguistic: its members are formulas, not sets inside a
model.  It is needed because transfer between transitive structures depends
on formula complexity.  In particular, for \(n\geq1\), \(\Delta_n\) is an
equivalence class designation relative to the base theory, not merely the
literal intersection of two strings; \(\Delta_0\) remains the syntactic
bounded class.  The L\'evy hierarchy is distinct from the projective
hierarchy of Definition~\ref{def:p06-projective-sigma12}.
The historical hierarchy source is \cite{Levy1965}; the modern
normal-form and absoluteness formulation used here follows \cite{Jech2003}.
\index{Levy hierarchy@L\'evy hierarchy}
\index{bounded quantifier}
\index{Delta-zero formula@\(\Delta_0\) formula}
\index{Sigma-n formula@\(\Sigma_n\) formula}
\index{Pi-n formula@\(\Pi_n\) formula}
\end{hbdefinition}

\begin{hbconvention}{def:p06-absoluteness}{Formula absoluteness}
Let \(M\subseteq N\) be transitive sets or classes, let
\(\vec a\in M^{<\omega}\), and let \(\varphi(\vec x)\) be a standard
membership-language formula.  The assertion that \(\varphi\) is
\emph{absolute between \(M\) and \(N\) at \(\vec a\)} means
\[
  (M,\in)\entails\varphi(\vec a)
  \quad\Leftrightarrow\quad
  (N,\in)\entails\varphi(\vec a),
\]
read formula by formula when either structure is a proper class.  This is a
metalinguistic relation among a formula, two membership structures, and a
parameter tuple.  It is not elementarity of \(M\) in \(N\): elementarity
requires agreement for every formula and every parameter tuple, whereas the
present notion may be restricted to one formula class.  The distinction is
used in the next theorem and in the forcing application below.
\index{absoluteness!set theory}
\end{hbconvention}

\begin{hbtheorem}{thm:p06-levy-absoluteness}{cited=Jech2003}{Bounded and one-quantifier absoluteness}
Let \(M\subseteq N\) be transitive sets or classes and let every displayed
parameter belong to \(M\).
\begin{enumerate}
  \item Every \(\Delta_0\) formula is absolute between \(M\) and \(N\).
  \item Every \(\Sigma_1\) formula is upward absolute: truth in \(M\)
        implies truth in \(N\).
  \item Every \(\Pi_1\) formula is downward absolute: truth in \(N\)
        implies truth in \(M\).
\end{enumerate}
In clauses 2 and 3, \(\Sigma_1\) and \(\Pi_1\) mean the displayed syntactic
normal-form representatives.  A formula classified only through a
base-theory equivalence inherits the corresponding transfer direction only
when that equivalence holds in both \(M\) and \(N\), as it does when both
structures satisfy the declared base.  Bare transitivity cannot import an
equivalence available only as a theorem of that base.
Transitivity and parameter membership are load-bearing.  No corresponding
unqualified rule is asserted for arbitrary, possibly nonstandard, models.
The forcing-specialized source sketch remains separately recorded in
Section~\ref{sec:constructibility-forcing-02}; this cited general statement
does not upgrade that sketch.
\end{hbtheorem}

\begingroup\small
\begin{longtable}{@{}p{.16\textwidth}p{.29\textwidth}p{.20\textwidth}p{.25\textwidth}@{}}
\caption{L\'evy levels and their elementary transfer behavior between
transitive \(M\subseteq N\) containing all parameters.}
\label{tab:p06-levy-absoluteness}\\
\toprule
Class & Normal-form description & Guaranteed direction & What is not licensed \\
\midrule
\endfirsthead
\toprule
Class & Normal-form description & Guaranteed direction & What is not licensed \\
\midrule
\endhead
\(\Delta_0\) & Every quantifier is bounded by a displayed set & Both
directions & Elementarity for formulas with unbounded quantifiers \\
\(\Sigma_1\) & Unbounded existential block over a bounded matrix & Upward,
from \(M\) to \(N\), for the normal form & Downward transfer, or use of a
base-only equivalence in structures that need not satisfy the base \\
\(\Pi_1\) & Unbounded universal block over a bounded matrix & Downward,
from \(N\) to \(M\), for the normal form & Upward transfer, or use of a
base-only equivalence in structures that need not satisfy the base \\
\(\Sigma_n,\Pi_n\), \(n\geq2\) & Alternating unbounded blocks, modulo
provable equivalence & No bare transitivity theorem & Projective or forcing
absoluteness without its additional hypotheses \\
\bottomrule
\end{longtable}
\endgroup

\section{Reflection and Shoenfield absoluteness}
\label{sec:p06-reflection-absoluteness}

\begin{hbtheorem}{thm:p06-levy-reflection}{cited=Jech2003}{Finite-family L\'evy reflection}
Assume \(\mathsf{ZF}\).  Let \(\Phi\) be a finite set of standard formulas in
the language of membership and let \(\beta\) be an ordinal.  There is an
ordinal \(\alpha>\beta\) such that, for every
\(\varphi(x_1,\ldots,x_m)\in\Phi\) and every
\(a_1,\ldots,a_m\in V_\alpha\),
\[
  \varphi(a_1,\ldots,a_m)
  \quad\Leftrightarrow\quad
  (V_\alpha,\in)\entails
    \varphi(a_1,\ldots,a_m).
\]
The left side is the ambient assertion for the one displayed standard
formula; it is not an application of a class satisfaction predicate to
\(V\).
The historical reflection source is \cite{Levy1960}.
\end{hbtheorem}

\begin{proofboundary}{def:p06-reflection-boundary}{No all-formula rank}
Theorem~\ref{thm:p06-levy-reflection} reflects each fixed finite family at
arbitrarily high ranks.  It does not produce one rank-initial segment that is
elementary in the universe simultaneously for all formulas.  That overclaim would
also ignore the theorem's metalinguistic, schema-level standing.
\end{proofboundary}

\begin{hbdefinition}{def:p06-projective-sigma12}{Boldface \(\boldsymbol\Sigma^1_2\) and \(\boldsymbol\Pi^1_2\) sentences}
With a real parameter \(a\), a boldface \(\boldsymbol\Sigma^1_2(a)\) sentence
is a second-order-arithmetic sentence equivalent to one of the form
\[
  \exists x\in\omega^\omega\,
  \forall y\in\omega^\omega\,
  \psi(x,y,a),
\]
where \(\psi\) is arithmetical; a
\(\boldsymbol\Pi^1_2(a)\) sentence has the dual universal--existential form.
Here \(\omega^\omega\) denotes the set of functions from \(\omega\) to
\(\omega\), not ordinal exponentiation; Chapter~6.7 equips that set with the
topology of Baire space.
These are metalinguistic projective-complexity classes of statements about
reals.  The boldface designation permits the displayed real parameter, unlike
the parameter-free lightface hierarchy, and it is needed to state exactly the
parameter range of Shoenfield absoluteness.  It is not a level of the L\'evy
hierarchy, whose unbounded quantifiers range over arbitrary sets.
\index{projective hierarchy}
\index{Shoenfield absoluteness}
\index{Sigma-one-two@\(\boldsymbol\Sigma^1_2\)}
\index{Pi-one-two@\(\boldsymbol\Pi^1_2\)}
\end{hbdefinition}

\begin{hbtheorem}{thm:p06-shoenfield-absoluteness}{cited=Moschovakis2009}{Shoenfield absoluteness}
Let \(M\subseteq N\) be transitive models of \(\mathsf{ZFC}\) with the same
ordinals, and let the real parameter \(a\) belong to \(M\).  Then \(M\) and
\(N\) agree on every boldface \(\boldsymbol\Sigma^1_2(a)\) sentence and,
by complementation, on every boldface \(\boldsymbol\Pi^1_2(a)\) sentence.
The hypotheses do not reduce to agreement about \(\omega\), and the theorem
does not assert absoluteness for all higher projective levels.
The original theorem is \cite{Shoenfield1961}; the displayed modern
second-order formulation is supported by \cite{Moschovakis2009}.
\end{hbtheorem}

\begin{hbexample}{def:p06-shoenfield-forcing}{Ground-model projective statements in a forcing extension}
Let \(M\) be the transitive model fixed by the convention in
Section~\ref{sec:constructibility-forcing-02}, let \(G\) be \(M\)-generic, and let
\(a\in M\) be a real.  The cited forcing theorem gives \(M\) and \(M[G]\) the
same ordinals.  Hence Theorem~\ref{thm:p06-shoenfield-absoluteness} gives
agreement on every
\(\boldsymbol\Sigma^1_2(a)\) or \(\boldsymbol\Pi^1_2(a)\) sentence.  This
application combines two cited results; the forcing diagram alone proves
neither same ordinals nor projective absoluteness
\cite{Jech2003,Moschovakis2009}.
\end{hbexample}

\section{Constructibility interface}
\label{sec:p06-constructibility-interface}

The constructible hierarchy and its formula-by-formula relativization are
defined in Section~\ref{sec:constructibility-forcing-01}.  The exact cited
constructibility scheme there includes \(\mathsf{GCH}\) in \(L\), and
Theorem~\ref{thm:p06-condensation} states the exact elementary-substructure
hypothesis under which a collapse is an earlier \(L_\beta\).  Jensen fine
structure remains a further-reading pointer, not a local entry.  These facts,
together with finite-family reflection above, fill the constructibility and
reflection cluster without identifying three different mechanisms:
condensation describes suitable elementary substructures of levels of \(L\),
reflection compares finitely many formulas with a rank \(V_\alpha\), and
absoluteness transfers restricted formula classes between transitive models
\cite[Chapters~12--13]{Jech2003}.  For condensation and the fine-structure
boundary, see \cite{Devlin1984}; no fine-structure theorem is imported here.

\section{The conditional large-cardinal ladder}
\label{sec:p06-large-cardinal-ladder}

Throughout this section, \emph{inaccessible} means strongly inaccessible, and
every large-cardinal assertion is conditional.  No theorem below says that
\(\mathsf{ZFC}\) proves the existence of a cardinal on the ladder.

\begin{hbdefinition}{def:p06-inaccessible-cardinal}{Strongly inaccessible cardinal}
An uncountable cardinal \(\kappa\) is \emph{strongly inaccessible} iff it is
regular and a strong limit:
\[
  \operatorname{cf}(\kappa)=\kappa
  \quad\text{and}\quad
  2^\lambda<\kappa\quad\text{for every cardinal }\lambda<\kappa.
\]
In this chapter the unqualified word \emph{inaccessible} has exactly this
meaning.  It names a property of a cardinal, not an existence axiom; the
regularity and strong-limit clauses distinguish it from merely uncountable or
weakly inaccessible cardinals.  The distinction is used in the hierarchy
model theorem and at the bottom of Table~\ref{tab:p06-large-cardinal-ladder}.
\index{inaccessible cardinal}
\index{strong-limit cardinal}
\end{hbdefinition}

\begin{hbdefinition}{def:p06-mahlo-cardinal}{Mahlo cardinal}
Under the convention used here, a cardinal \(\kappa\) is \emph{Mahlo} iff it
is strongly inaccessible and the regular cardinals below \(\kappa\) form a
stationary subset of \(\kappa\), in the sense of
Definition~\ref{def:p06-club-stationary}.  Equivalently under this convention,
the strongly inaccessible cardinals below \(\kappa\) form a stationary set.
This is a cardinal property strengthening inaccessibility by internal
reflection of regularity.  Authors vary between weak and strong Mahlo
conventions, so the displayed clause, rather than the bare name, controls all
later uses in this handbook.
\index{Mahlo cardinal}
\end{hbdefinition}

\begin{hbdefinition}{def:p06-weakly-compact-cardinal}{Weakly compact cardinal}
An uncountable cardinal \(\kappa\) is \emph{weakly compact} iff
\[
  \kappa\longrightarrow(\kappa)^2_2
\]
in the partition notation of Definition~\ref{def:p06-partition-arrow}: every
two-coloring of \([\kappa]^2\) has a homogeneous subset of cardinality
\(\kappa\).  This is a partition property of a cardinal, not topological
compactness and not a rule of implication.  It is needed for the combinatorial
calibration below and is stronger than the tree property stated without the
inaccessibility hypothesis.
\index{weakly compact cardinal}
\end{hbdefinition}

\begin{hbdefinition}{def:p06-measurable-cardinal}{Measurable cardinal}
An uncountable cardinal \(\kappa\) is \emph{measurable} iff some nonprincipal
\(\kappa\)-complete ultrafilter \(\mathcal U\) exists on \(\kappa\).  Here
\(\kappa\)-complete means that the intersection of every family of fewer than
\(\kappa\) members of \(\mathcal U\) lies in \(\mathcal U\), and nonprincipal
means that no singleton belongs to \(\mathcal U\).  The term names a cardinal
property witnessed by
an ultrafilter; it is not ordinary measure-theoretic measurability.  The
completeness clause is precisely what supports the well-founded ultrapower in
Theorem~\ref{thm:p06-measurable-ultrapower}.
The measure characterization originates with \cite{Ulam1930}.
\index{measurable cardinal}
\index{ultrafilter!kappa-complete@\(\kappa\)-complete}
\end{hbdefinition}

\begin{hbconvention}{def:p06-critical-point}{Critical point and closure notation}
For a nontrivial elementary embedding \(j:V\to M\) into a transitive class,
\(\operatorname{crit}(j)\) denotes the least ordinal \(\alpha\) such that
\(j(\alpha)\ne\alpha\).  The assertion \({}^{\lambda}M\subseteq M\) means
that every function \(f:\lambda\to M\) belonging to the ambient universe is
itself in \(M\).  Both are metalinguistic class-embedding conventions, read
formula by formula; they do not add a class object or satisfaction predicate
to \(\mathsf{ZFC}\).  The closure notation is stronger than transitivity and
is introduced because it is the load-bearing extra condition in the next
definition.
\index{critical point!elementary embedding}
\index{elementary embedding!set theory}
\end{hbconvention}

\begin{hbdefinition}{def:p06-supercompact-cardinal}{Supercompact cardinal}
A cardinal \(\kappa\) is \emph{supercompact} iff, for every cardinal
\(\lambda\geq\kappa\), there are a transitive class \(M\) and a nontrivial
elementary embedding \(j:V\to M\) such that
\[
  \operatorname{crit}(j)=\kappa,
  \qquad j(\kappa)>\lambda,
  \qquad {}^{\lambda}M\subseteq M.
\]
This is an embedding characterization of a cardinal property.  Unlike
measurability, it requires the displayed closure at every target
\(\lambda\geq\kappa\); that distinction places supercompactness above
measurability in the one-way implication ladder.  No strict consistency
separation is asserted.  The alternative fine-normal-measure definition is
not mixed into this entry.
\index{supercompact cardinal}
\end{hbdefinition}

\begin{hbtheorem}{thm:p06-inaccessible-stage}{cited=Kanamori2003}{An inaccessible hierarchy stage models ZFC}
If \(\kappa\) is strongly inaccessible, then the set-sized membership
structure \((V_\kappa,\in)\) satisfies \(\mathsf{ZFC}\).  This is the
conditional annotation on the exact rank axis in
Figure~\ref{fig:p06-cumulative-hierarchy}; the figure asserts no inaccessible
cardinal exists.
\end{hbtheorem}

\begin{hbtheorem}{thm:p06-weak-compact-tree}{cited=Kanamori2003}{Tree characterization of weak compactness}
For an uncountable cardinal \(\kappa\), the following are equivalent:
\begin{enumerate}
  \item \(\kappa\longrightarrow(\kappa)^2_2\);
  \item \(\kappa\) is strongly inaccessible and every tree of height
        \(\kappa\), all of whose levels have cardinality below \(\kappa\),
        has a branch meeting every level.
\end{enumerate}
Thus the tree property without inaccessibility is not the stated
characterization.
\end{hbtheorem}

\begin{hbtheorem}{thm:p06-measurable-ultrapower}{cited=Kanamori2003}{Measurability and the universe ultrapower}
If \(\mathcal U\) is a nonprincipal \(\kappa\)-complete ultrafilter on
\(\kappa\), the class ultrapower \(\operatorname{Ult}(V,\mathcal U)\) is well
founded.  Its transitive
collapse yields a nontrivial elementary embedding
\[
  j_{\mathcal U}:V\longrightarrow M
  \qquad\text{with}\qquad
  \operatorname{crit}(j_{\mathcal U})=\kappa.
\]
Conversely, such an embedding with critical point \(\kappa\) recovers a
nonprincipal \(\kappa\)-complete ultrafilter by
\[
  \mathcal U=\{X\subseteq\kappa:\kappa\in j(X)\}.
\]
The complete class-ultrapower construction is cited; this statement is not a
compressed local proof of its coding and well-foundedness obligations.
The ultrapower role is historically attached to \cite{Keisler1962}; the
modern large-cardinal formulation is \cite{Kanamori2003}.
\end{hbtheorem}

\begin{hbtheorem}{thm:p06-large-cardinal-implications}{cited=Kanamori2003}{Implications through supercompactness}
For the definitions fixed above,
\[
  \begin{gathered}
  \text{supercompact}
    \quad\text{implies}\quad \text{measurable}
    \quad\text{implies}\quad \text{weakly compact},\\
  \text{weakly compact}
    \quad\text{implies}\quad \text{Mahlo}
    \quad\text{implies}\quad \text{strongly inaccessible}.
  \end{gathered}
\]
These are implications between conditional existence properties.  The display
does not assert that \(\mathsf{ZFC}\) proves any rung exists, that adjacent
consistency strengths are strictly separated, or that the chart exhausts the
large-cardinal hierarchy.
\end{hbtheorem}

\begin{hbtheorem}{thm:p06-scott-measurable}{cited=Scott1961}{Scott's constructibility obstruction}
If a measurable cardinal exists, then \(V\ne L\).  The theorem does not say
that every large-cardinal property is incompatible with constructibility.
\end{hbtheorem}

For the calibration in the next table, inaccessible, Mahlo, and weakly
compact cardinals can occur in \(L\), whereas supercompactness implies the
measurability already excluded by Scott's theorem.  Thus measurable is the
first displayed rung incompatible with \(V=L\).  Here ``can occur in \(L\)''
is a relative compatibility claim conditional on the corresponding
large-cardinal hypothesis, not an existence assertion \cite{Kanamori2003}.

\begingroup\scriptsize
\begin{longtable}{@{}p{.15\textwidth}p{.28\textwidth}p{.27\textwidth}p{.20\textwidth}@{}}
\caption{Large-cardinal ladder.  Every row is conditional;
rows run from weaker to stronger, and implications run upward from a lower
row.  This is a consistency-strength calibration, not a ZFC existence theorem
or a proof of strict separation.}
\label{tab:p06-large-cardinal-ladder}\\
\toprule
Rung & Defining property & Exact calibration used here & Relation to \(V=L\) \\
\midrule
\endfirsthead
\toprule
Rung & Defining property & Exact calibration used here & Relation to \(V=L\) \\
\midrule
\endhead
Strongly inaccessible & Uncountable, regular, and strong limit &
\((V_\kappa,\in)\entails\mathsf{ZFC}\) & Can occur in \(L\) \\
Mahlo & Inaccessible; regular cardinals below \(\kappa\) are stationary &
Reflects regularity stationarily often below \(\kappa\) & Can occur in \(L\) \\
Weakly compact & \(\kappa\longrightarrow(\kappa)^2_2\) & Equivalently,
inaccessible plus the displayed \(\kappa\)-tree branch property & Can occur in
\(L\) \\
Measurable & Carries a nonprincipal \(\kappa\)-complete ultrafilter &
Well-founded universe ultrapower with critical point \(\kappa\) & Impossible
under \(V=L\) by Scott's theorem \\
Supercompact & For every \(\lambda\geq\kappa\), an embedding with
\(j(\kappa)>\lambda\) and \({}^{\lambda}M\subseteq M\) & Implies measurable;
the table stops here rather than claiming an exhaustive hierarchy & Impossible
under \(V=L\), since it implies measurable \\
\bottomrule
\end{longtable}
\endgroup

Stronger large-cardinal hypotheses continue beyond the displayed
supercompact rung; \cite{Kanamori2003} is the reader pointer.  No additional
large-cardinal term or implication is introduced locally because the register
ceiling ends at supercompactness.

\begin{hbwarning}{def:p06-ctm-inaccessible-boundary}{An inaccessible stage is not equivalent to an arbitrary ctm}
If an inaccessible \(\kappa\) exists, then
Theorem~\ref{thm:p06-inaccessible-stage}, followed externally by downward
L\"owenheim--Skolem and Mostowski collapse, yields a countable transitive
model of \(\mathsf{ZFC}\).  The converse is false: an arbitrary transitive set
model of \(\mathsf{ZFC}\) need not be an ambient rank \(V_\kappa\), and its
ordinal height need not be an inaccessible cardinal of the ambient universe
\cite{Jech2003,Kanamori2003}.
Accordingly, the forcing convention and its countable-transitive-set strength
analysis in Section~\ref{sec:constructibility-forcing-02} are not replaced by
a large-cardinal existence assumption.
\end{hbwarning}

%% Source/status disposition for verification:
%% - no v1.1 proof or sketch is restated with a stronger status;
%% - all theorem-class additions in this file are cited;
%% - L/GCH/condensation remain in constructibility-forcing.tex;
%% - A-130 remains fig:p06-cumulative-hierarchy;
%% - A-069/A-134 are the following consistency-strength file.
